% ─────────────────────────────────────────────────────────────────
\section{Data Pipeline: VTK Mesh to SSL Graph}
\label{sec:pipeline}
% ─────────────────────────────────────────────────────────────────

The raw SHREC\,25 data consists of VTK surface meshes with
$\approx$10\,676 vertices.
Two preprocessing stages convert each mesh into a compact graph suitable for
message-passing GNNs.

\subsection{Stage 1 --- SSL Graph Construction}
\label{sec:pipeline:ssl}

Implemented in \texttt{vtk\_mesh\_to\_ssl\_graph\_parallel.py}.

\paragraph{Farthest-point subsampling (FPS).}
Each mesh is subsampled to exactly $N_v = 1\,000$ nodes using farthest-point
sampling in Euclidean 3D space.
FPS produces a spatially uniform coverage of the surface, avoiding the
concentration artefacts of random subsampling.

\paragraph{k-NN graph construction.}
A $k$-nearest-neighbour graph with $k=6$ is built over the 1\,000 FPS
nodes.
Edges are undirected; we retain no edge weights (all message-passing
operations treat the graph as unweighted).

\paragraph{Feature extraction.}
For each retained node the following 8-dimensional feature vector is computed:
\begin{equation}
  \mathbf{x}_i = \bigl[
    \underbrace{\hat{P}_i}_{\text{Potential}},\;
    \underbrace{\hat{N}_i}_{\text{NormPot}},\;
    \underbrace{n_x, n_y, n_z}_{\text{normals}},\;
    \underbrace{x_\text{ssl}, y_\text{ssl}, z_\text{ssl}}_{\text{coords}}
  \bigr],
  \label{eq:feature_vec}
\end{equation}
where $\hat{P}_i$ and $\hat{N}_i$ are the z-scored, $[\!-5,5\!]$-clamped
versions of \textit{Potential} and \textit{NormalPotential}, and
$(x_\text{ssl}, y_\text{ssl}, z_\text{ssl})$ are the node coordinates
translated so that their centroid is at the origin and scaled to unit
maximum absolute value.
Normals $(n_x, n_y, n_z)$ are taken directly from the VTK file and
renormalised to unit length.

The VTK reader must be invoked with \texttt{ReadAllScalarsOn()} to load both
scalar arrays; otherwise only the default (first) array is available.

\paragraph{Output format.}
Each graph is saved as a PyTorch Geometric \texttt{Data} object with fields
\texttt{x} (float32, $1000 \times 8$), \texttt{edge\_index} (long, $2 \times
E$), and \texttt{y} (int64 class label).
Separate \texttt{train\_split.csv} and \texttt{val\_split.csv} are generated
with an 80/20 stratified split.

\subsection{Stage 2 --- Geodesic Patch Decomposition}
\label{sec:pipeline:patch}

Implemented in \texttt{build\_patchassign\_dataset.py}.

\paragraph{Patch seed selection.}
$k$ seed nodes are selected by another round of FPS over the 1\,000-node
graph, producing spatially well-distributed patch centres.
We experiment with $k \in \{8, 32, 64, 128\}$; \textbf{k=64} is used in all
best-performing runs.

\paragraph{Geodesic Voronoi assignment.}
Each non-seed node is assigned to the patch whose centre is geodesically
nearest, using Euclidean edge lengths as a proxy for geodesic distance
(computed via Dijkstra's algorithm over the k-NN graph).
Isolated nodes that cannot be reached from any seed are assigned to the
nearest centre by Euclidean fallback.

\paragraph{Patch meta-graph.}
A patch-level meta-graph is constructed in which each patch is a node and
two patches are connected if they contain adjacent nodes in the original graph.
Patch-level features are the mean of their constituent node features (used only
as initial features before the patch GIN updates them).

\paragraph{Output format.}
For each graph a companion \texttt{.pt} file stores
\texttt{patch\_ids} (int64, $1000$-vector mapping nodes to patch indices),
\texttt{patch\_edge\_index} (long, $2 \times E_p$), and
\texttt{num\_patches} (int).

\subsection{Summary}

The two-stage pipeline reduces each mesh from $\sim$10\,000 vertices to 1\,000
graph nodes, then to 64 patch nodes.
This represents a $10\times$ and $156\times$ compression of the original mesh,
respectively, while preserving the spatial and electrostatic structure needed
for classification.
The end-to-end pipeline from raw VTK to classification-ready tensors processes
all 11\,555 proteins in under 30 minutes on a single CPU core.
